\chapter{Proofs for Morley's Original Categoricity Argument}
\label{annex:p03-n-a-006-01}

\begin{annexitem}{annex:p3.morley.normalization}{thm:p03-morley-normalization}{The original normalization and amalgamation setup}
\begin{proof}
The initial graph translation is truth preserving. For each term,
inductively define its value by existentially naming the values of
its immediate subterms and using the appropriate function graph.
Totality and single-valuedness give one value, including the
singleton graph for a constant. Translate atomic formulas using
these term values; Boolean and quantifier induction then preserves
all formulas. Conversely the graphs recover unique functions and
constants on every model. This preserves model cardinalities,
isomorphisms and elementary maps and permits arbitrary induced
parameter subsets in the resulting purely relational language.
For each original formula add the axiom identifying its new relation
with that formula. These axioms force its interpretation and give
the unique expansion of every original model. Translate an expanded
formula back by replacing new relation symbols with their definitions,
renaming bound variables to avoid capture. The original formula
obtained has its own relation symbol, giving atomic equivalence in
the expansion. Completeness, isomorphisms, elementary maps and
cardinalities therefore transfer. Realization of parameter types
transfers by the same formula translation, so saturation transfers.

The named diagram of an induced subsystem is its entire elementary
parameter information in the original language. If two extensions
of a common subsystem \(A\) are given, their complete elementary
diagrams can be united, identifying only the names of \(A\). This
union is consistent. Indeed, an inconsistent finite subset would
give formulas \(\theta_1(\bar a,\bar b_1)\) and
\(\theta_2(\bar a,\bar b_2)\) whose conjunction is impossible;
the two existential projections in the variables \(\bar b_i\)
are both true in the same complete parameter theory over \(A\),
and that theory entails their conjunction. A model of the union
is an amalgam. The inequalities inside each elementary diagram
make both embeddings injective; no disjointness between their
unnamed parts is asserted. Amalgamating over the empty parameter
system gives joint embedding.

For an increasing family of subsystems, each finite tuple occurs
in one stage. Its complete elementary diagram is therefore
consistent with \(T\); compactness embeds the union as an induced
subsystem of a model. If the stages are themselves models, their
inclusions are elementary, and the elementary-chain theorem
makes the union a model. Alternatively, every existential
formula true over a tuple from a stage has a witness in that stage.
This also proves the needed chain theorem directly. These are
the contents of the source's Lemmas 1.1--1.2.
\end{proof}
\end{annexitem}

\begin{annexitem}{annex:p3.morley.ramsey-em}{thm:p03-morley-ramsey-em}{The finite-color argument and the original term hulls}
\begin{proof}
For Ramsey's assertion reduce to a countably infinite subset.
Induct on \(n\). The case \(n=1\) is the infinite pigeonhole
principle; \(n=0\) is immediate. For a coloring of
\((n+1)\)-subsets, recursively choose distinct \(x_j\) and
nested infinite remaining sets. After choosing \(x_j\),
apply the induction hypothesis to the coloring of
\(n\)-subsets of the remaining set obtained by adjoining
\(x_j\). Retain an infinite homogeneous remaining set
and call its color \(c_j\). Infinitely many \(j\)'s
have the same \(c_j=c\). The corresponding \(x_j\)'s
are homogeneous: in any increasing \((n+1)\)-tuple,
the last \(n\) choices lie in the remaining set chosen
after its first member, so its color is \(c\).

Now add a distinct constant for each \(x\in X\),
pairwise inequalities, and the equivalences asserting
that every two increasing tuples of the same length
agree on every formula. Each finite subset mentions
finitely many formulas and finitely many constants.
In an infinite original model, color the finite tuples
by the finitely many truth values required. Repeated
applications of Ramsey, one for each of the finitely
many arities, leave an infinite set homogeneous for
all these truth vectors. Interpret the finitely many
constants in its order, preserving their finite
ordering pattern. Thus every finite constraint is
satisfiable, and compactness gives the required
ordered indiscernibles. This is the source's Theorem 3.2.

Here is also the precise Skolem construction used by
Theorem 3.3. Put each axiom in prenex form; replace
each existentially quantified variable by a new
function of the universally quantified variables
preceding it, or a constant if there are none.
The resulting axioms are universal and have quantifier-free
matrices. The language remains countable. Every
original model expands to one satisfying them by
choosing the successive witnesses, and every reduct
of a model of them satisfies the original prenex
axioms. Hence the expansion is conservative for
original sentences. Include a fresh constant if
needed so that the generated hull is nonempty.
Substructures closed under these functions satisfy
the universal theory, which proves that the term
hull is a model of it.

Apply the ordered-indiscernible construction to this
expanded theory. A tuple of term values is governed
by the ordered finite union of its generators.
Consequently replacing those generators by an
order-isomorphic finite tuple preserves every
formula on the term values. In particular it
preserves their equalities; the assignment
\[
 t(x_1,\ldots,x_r)\longmapsto
           t(f(x_1),\ldots,f(x_r))
\]
is well defined even if an element has different
term representations. It preserves all relations,
functions and inequalities, and is onto when
\(f\) is onto. This gives Theorems 3.4--3.5.
Performing the construction on \(\Sigma^{*\#}\)
makes the hull reduct a \(\Sigma^*\)-model and its
inclusion elementary on the original language,
by normalization. The induced generator maps are
therefore elementary on the original reduct,
giving both clauses of Theorem 3.6. Countably many
finite-arity terms on an infinite set of size
\(\kappa\) have at most \(\kappa\) values, while
the distinct generators give at least \(\kappa\).
\end{proof}
\end{annexitem}

\begin{annexitem}{annex:p3.morley.thin-models}{thm:p03-morley-thin-models}{Ordinal generators and the contradiction to categoricity}
\begin{proof}
Use the initial ordinal \(\kappa\) as the ordering
of the generators in the Morleyized Skolem hull.
Let \(Y\) be a countable subset of its original
reduct. Choose a term representation of each
\(y\in Y\); the union \(X_0\) of their finite
supports is countable. Every other element is
a value of one of countably many terms on a
finite tuple of generators. Its type over \(Y\)
is determined by that term and the ordering
of its generator tuple relative to \(X_0\):
an order isomorphism fixing \(X_0\) carries
every named \(y\) to itself and preserves all
formulas. A countable subset of an ordinal
has a countable well-order type. It has only
countably many cuts, so finitely many new
generators have only countably many such
interpolation patterns, including repetitions
and mutual order. The model therefore realizes
only countably many types over \(Y\).

If the theory is not totally transcendental,
the residual-tree theorem gives a countable
\(C\) with uncountably many types. Select
\(\aleph_1\) different types. Simultaneously
realize them in an extension, then apply
Löwenheim--Skolem and, if needed, its upward
form to obtain a model of the categorical
uncountable size \(\kappa\) containing
\(C\) and all the selected realizations.
Their number is at most \(\kappa\). This
model has a countable subset over which it
realizes uncountably many types, while the
thin model does not have any such subset.
The property is invariant under isomorphism,
contradicting categoricity. This proves
Theorems 3.7--3.8 without assuming that
\(\kappa\) can contain all continuum many
types simultaneously.

For dense linear order, different Dedekind
cuts of the rationals give consistent
different one-types, distinguished by a
named rational between the cuts. There are
\(2^{\aleph_0}\) cuts, although all countable
dense orders without endpoints are isomorphic
by the elementary back-and-forth construction.
Thus the countable version of the hypothesis
would not suffice.
\end{proof}
\end{annexitem}

\begin{annexitem}{annex:p3.morley.no-order}{thm:p03-morley-no-order}{The permutation pattern creates continuum many types}
\begin{proof}
Suppose an \(n\)-ary relation has both stated
properties on an infinite set. For \(n<2\)
there is only one relevant permutation,
so both properties cannot hold. For
\(n\geq2\), impose an arbitrary order
on that set and color each increasing
\(n\)-tuple by its truth vector under all
\(n!\) permutations. Ramsey gives an
infinite subset with a constant truth vector.
The vector has both a true and a false
entry. Adjacent transpositions generate
the permutation group, so some pair of
entries differing by one adjacent exchange
has different truth values.

Compactness reproduces this same vector
on an ordered family of any desired order
type: a finite set of constraints is
realized in the original homogeneous
infinite family. Use order type
\(\mathbb R\), and let \(Z\) be its
countable rationally ordered dense subset.
For any \(y<y'\) outside \(Z\), put
one point of \(Z\) between them and
choose the other \(n-2\) points on either
side so that \(y\) and this middle point
occupy the adjacent ranks at which the
two chosen permutation entries differ.
Substituting \(y'\) for \(y\) exchanges
these ranks and leaves the other ranks
unchanged. The resulting relation formula,
with its \(n-1\) named parameters in
\(Z\), has different truth values at
\(y\) and \(y'\). Thus different real
points outside \(Z\) have different
types over \(Z\), giving continuum
many types over a countable subsystem.
This contradicts total transcendence.

Naming finitely many parameters does not
change the countability characterization:
one-types over a countable set in the
named theory are original types over
that set together with those parameters.
Thus the named complete theory is also
totally transcendental, and the preceding
argument applies to any definable relation
with finite parameters. This is Theorem 3.9.
\end{proof}
\end{annexitem}

\begin{annexitem}{annex:p3.morley.prime-models}{thm:p03-morley-prime-models}{Density, isolation, and the prime-extension construction}
\begin{proof}
Embed \(A\) in a model realizing all its types.
Tarski--Vaught says that \(A\) is elementary
in that model exactly when every nonempty
one-variable definable set over \(A\)
has a witness in \(A\). These sets are
exactly the nonempty basic clopens, so
the criterion is density of its realized
types. This proves Lemma 4.1.

In a totally transcendental theory choose
a point of least rank \(\alpha\) in a
nonempty open \(U\). Every point of
\(U\) belongs to \(S^\alpha(A)\).
The chosen point has an isolating
neighborhood in that derivative; intersect
it with \(U\) to isolate it in \(S(A)\)
itself. This proves Lemma 4.2.

For the more general existence assertion,
write \(\mu=\kappa(A)+\aleph_0\).
There are at most \(\mu\) isolated
points, because each has a distinct
isolating basic clopen. List them.
Starting at \(A\), deal with each
listed point: if already realized,
do nothing; otherwise its inverse
image over the current subsystem is
a nonempty clopen. Density there
provides an isolated lift. Adjoin
one realizing element. At limits
take unions of subsystems.

Every map of the starting system
into a model extends through this
construction. For a new isolated
type, translate its isolating formula
along the current map. It is consistent
over the image parameter system, so
the target model contains a witness.
All formulas of the isolated type
then hold at that witness. If the
new point was outside the current
system its inequalities from old
elements prevent a collision with
their images. Hence the map extends
as a monomorphism.

After \(\mu\) steps the resulting
system realizes all isolated types
over \(A\). Repeat this entire sweep
over the resulting system, and
continue for \(\omega\) sweeps.
Each sweep has at most \(\mu\)
steps and the size stays at most
\(\mu\). Let \(B\) be the union.
Every parameter formula over \(B\)
already occurs over one sweep's
starting system. If it is nonempty,
density of isolated points there
gives an isolated type containing it,
realized by the next sweep. Thus
the realized points are dense in
\(S(B)\), and \(B\) is a model.
The stagewise extension of maps
makes it prime over \(A\).
This is the full construction of
Theorem 4.3.
\end{proof}
\end{annexitem}

\begin{annexitem}{annex:p3.morley.chain-types}{thm:p03-morley-chain-types}{The isolation argument at a limit stage}
\begin{proof}
Along a coherent family the ordinal
ranks are nonincreasing by exact
images. They eventually stabilize,
since a strictly decreasing infinite
sequence of ordinals is impossible.
Once the rank has stabilized, the
positive integer degrees are
nonincreasing by addition and hence
stabilize too. Call the eventual
values \(\nu,n\). The union of
the types is consistent by finite
character and complete because each
parameter formula occurs at a stage.
It is consequently the unique union
type \(p\).

For every sufficiently late index
\(\xi\), the compact fiber over
\(p_\xi\) in \(S^\nu(A)\) is
nonempty. These fibers are nested;
their intersection is nonempty and
can only be \(\{p\}\). Thus
\(p\) has rank at least \(\nu\),
and restriction to \(p_\xi\)
bounds it above by \(\nu\).
At any later stage \(\eta\), the
point \(p_\eta\) already has degree
\(n=d(p_\xi)\) and is the sole
rank-\(\nu\) lift of \(p_\xi\).
Every rank-\(\nu\) point in the
fiber over \(p_\xi\) in \(S(A)\)
must therefore restrict to every
\(p_\eta\), and hence equal \(p\).
Addition now gives \(d(p)=n\).
This proves part (a) of Lemma 4.4,
including its last assertion.

For part (b), insert union stages
if the original chain is not
continuous. Starting with the
isolated \(p_0\), at a successor
choose a minimum-rank point in its
nonempty clopen inverse image.
The density argument makes this
point isolated. At a limit the
previous coherent types have the
unique union type by part (a).
We must prove its isolation, not
merely its existence.

Choose a late stage \(\xi_0\)
where rank \(\nu\) and degree
\(n\) have stabilized. The
inverse image of the isolated
singleton \(\{p_{\xi_0}\}\)
is a clopen neighborhood of the
union type. Its rank-\(\nu\)
part is the singleton union
type, by part (a) and the
degree law. Suppose it had
another point \(q\) of lower
rank. Finite localization of
\(q\)'s rank supplies a stage
at which its restriction has
rank less than \(\nu\).
Let \(\eta\) be the first
such stage after \(\xi_0\).
Before \(\eta\) the restrictions
of \(q\) have rank \(\nu\),
and they must equal the
\(p_\xi\)'s: those coherent
points exhaust the degree
\(n\) of \(p_{\xi_0}\),
leaving no other rank-\(\nu\)
lift. The first drop cannot
be at a limit, because part
(a) would give rank \(\nu\)
for that union. At a successor
\(\eta=\zeta+1\), it contradicts
the choice of \(p_\eta\)
as minimum rank in the
fiber over \(p_\zeta\).
No such \(q\) exists, so
the clopen neighborhood is
the isolated singleton.
This also allows the final
union stage to be included.
\end{proof}
\end{annexitem}

\begin{annexitem}{annex:p3.morley.chain-primes}{thm:p03-morley-chain-primes}{Prime extensions that remain compatible with every later base}
\begin{proof}
Put \(A=\bigcup_{\xi<\gamma}A_\xi\).
Keep all \(B_\xi\)'s inside an
increasing family of ambient
subsystems containing \(A\).
The successor construction will
have the following stronger
property: for every later
\(s\geq\xi\), every map
\[
       A_s\cup B_{\xi-1}\longrightarrow D
\]
into a model extends to a map
\(A_s\cup B_\xi\to D\).
Here the predecessor notation
is used only at successors;
at \(\xi=0\) the starting
domain is \(A_s\). This
property controls the union
with the future base as well
as the prime extension itself.

At stage \(\xi=0\), or
\(\xi=\beta+1\), let the
starting subsystem be
\(K=A_\xi\) or
\(K=A_\xi\cup B_\beta\),
respectively. Run the isolated-type
sweeps of the prime construction,
but with an additional requirement.
When an isolated type over the
current \(K\) is to be realized,
apply the isolated-lift lemma
to the entire increasing chain
\[
  A_s\cup B_\beta\cup E
             \quad(\xi\leq s<\gamma),
\]
where \(E\) is the set of
elements already adjoined in
this sweep construction; omit
\(B_\beta\) at the initial
stage. The lemma provides
coherent isolated lifts over
every one of these later
systems, including their
union \(A\cup B_\beta\cup E\).
Realize that union type in an
extension of the ambient
subsystem. If it is already
realized there, use that
element instead. Either way,
the new element's type is
isolated over every later
base together with all
previously adjoined elements.

After the countably many
sweeps let \(B_\xi=K\cup E\).
The same witness-density
argument makes it a model.
For a fixed later \(A_s\),
each adjoined element has an
isolating formula over
\(A_s\cup B_\beta\) and
the earlier adjoined elements.
Thus every map of the starting
union into a model extends
element by element, exactly
as in the prime construction.
This proves the stronger
property. All choices are
consistent with the ambient
\(A\) because we realized
the union type, not a type
over the initial base alone.

At a limit \(\delta<\gamma\)
put \(B_\delta=\bigcup_{\xi<\delta}B_\xi\).
It is a model by the elementary
chain theorem, and contains
\(A_\delta\) by continuity.
To verify primality at any
index \(\delta\), begin with
a map \(A_\delta\to D\).
Use the stronger property at
each earlier successor, taking
the later base to be
\(A_\delta\), to extend
coherently to
\(A_\delta\cup B_\xi\).
At intervening limits take
the union of the maps. At
\(\delta\) this gives a map
of \(B_\delta\), or at a
successor apply that stage's
strong property once more.
Hence every \(B_\delta\)
is prime over \(A_\delta\).
This completes Theorem 4.5,
with the simultaneous later-base
isolation that its proof needs.
\end{proof}
\end{annexitem}

\begin{annexitem}{annex:p3.morley.indiscernible-extraction}{thm:p03-morley-indiscernible-extraction}{Minimal rank and degree produce a set, not just a sequence}
\begin{proof}
First suppose \(\kappa\) is
regular and uncountable. For
any \(C\) between \(A\)
and \(B\) of size less than
\(\kappa\), there are fewer
than \(\kappa\) types over
\(C\). Regularity implies
that some such type is
realized by \(\kappa\)
elements of \(B\). Among
all these pairs \((C,p)\)
choose \((C_0,p_0)\) of
least rank \(\nu\), and
then least degree \(n\)
at that rank.

Whenever \(C_0\subseteq C'\subseteq B\)
and \(\kappa(C')<\kappa\),
partition the \(\kappa\)
realizations of \(p_0\)
by their types over \(C'\).
Some extension \(p'\)
is still realized \(\kappa\)
times. Minimality and exact
images force its rank to
be \(\nu\); degree
minimality and the sum law
force its degree to be
\(n\). It is therefore
the unique rank-\(\nu\)
extension of \(p_0\).

Recursively choose distinct
\(x_\alpha\in B\setminus C_0\)
for \(\alpha<\kappa\),
with \(x_\alpha\) realizing
this unique extension over
\(C_\alpha=C_0\cup\{x_\beta:\beta<\alpha\}\).
The base has size less
than \(\kappa\), and
there are \(\kappa\)
choices before exclusion
of previously chosen points.
These extension types are
coherent and all have
rank \(\nu\), degree \(n\).

Increasing finite tuples
from the sequence have
the same type over
\(C_0\). Prove this by
induction on tuple length.
After mapping the first
\(m-1\) entries, the
last entry in each tuple
has rank \(\nu\), degree
\(n\) over those entries
and \(C_0\): it lies
between its type over
\(C_0\) and its type
over \(C_\alpha\), both
of those values. It is
the unique rank-\(\nu\)
extension of \(p_0\).
The already established
map carries one such
extension to the other,
so it extends to the
last entry. Thus this
is an order-indiscernible
sequence over \(C_0\).

It is in fact an
unordered indiscernible
set. If a formula with
finite parameters from
\(C_0\) depended on
the order of distinct
entries, its truth vector
on permutations of each
increasing tuple would
have both true and false
entries. It would define
a connected and antisymmetric
relation on this infinite
set. The no-order theorem,
applied after naming those
finite parameters, forbids
that. Hence all injective
maps fixing \(C_0\), and
therefore fixing \(A\),
preserve every finite
tuple formula.

For singular \(\kappa\)
and a given \(\lambda<\kappa\),
choose a regular successor
cardinal \(\rho<\kappa\)
larger than both \(\lambda\)
and \(\kappa(A)\).
Choose an induced subsystem
\(C\) of \(B\), containing
\(A\), of size \(\rho\),
and apply the regular case.
Such a \(\rho\) exists
because \(\kappa\) is a
limit cardinal; under
Choice, \(\mu^+\) is
regular since a union
of at most \(\mu\)
ordinals below \(\mu^+\)
has size at most \(\mu\).
The resulting set has
size \(\rho>\lambda\)
and lies inside the
given \(B\). This is
Theorem 4.6.
\end{proof}
\end{annexitem}

\begin{annexitem}{annex:p3.morley.saturation-existence}{thm:p03-morley-saturation-existence}{Back-and-forth and the length-\(\omega_1\) realization chain}
\begin{proof}
For uniqueness enumerate
two saturated models of
size \(\kappa\). Recursively
extend a partial elementary
bijection by one element
from the first enumeration
and then one from the
second. At every stage
\(\alpha<\kappa\) the
domain and range have
size less than \(\kappa\).
The next element's type
over the current domain,
translated by the map,
is realized by saturation
in the other model.
Inequalities from the
old domain ensure that
new elements have new
images. At limits take
unions. The resulting
total map is an elementary
bijection and hence an
isomorphism. This supplies
the uniqueness statement
that the source cites
as Theorem 5.1.

For the existence needed
in the proof, start with
any model \(B_0\) of
size \(\kappa>\aleph_0\).
There are at most
\(\kappa\) types over it
by total transcendence.
Realize all these in
an extension and retain
by Löwenheim--Skolem a
model \(B_1\) of size
\(\kappa\) containing
the old model and the
realizations. Repeat for
\(\omega_1\) stages,
taking elementary unions
at limits. Every stage
and the final union
have size \(\kappa\),
since \(\omega_1\leq\kappa\).
Any countable parameter
subsystem in the union
is contained in some
\(B_\alpha\), because
\(\omega_1\) is regular.
Every type over it extends
to one over \(B_\alpha\),
by onto restriction, and
is realized in
\(B_{\alpha+1}\). This
is exactly Theorem 5.2;
no cardinal exponentiation
condition was used.
\end{proof}
\end{annexitem}

\begin{annexitem}{annex:p3.morley.countable-omission-witness}{thm:p03-morley-countable-omission-witness}{Reducing an omitted type while retaining indiscernibles}
\begin{proof}
There are \(C\subseteq B\)
of size less than
\(\kappa(B)\) and
\(p\in S(C)\) omitted
in \(B\). The extraction
theorem provides an
infinite set \(Y\subseteq
B\setminus C\) indiscernible
over \(C\); retain a
countably infinite subset.
Take a countable elementary
submodel \(A_0\subseteq B\)
containing \(Y\).

For each \(a\in A_0\),
its type over \(C\)
differs from \(p\).
Choose a formula true
of \(a\) whose negation
belongs to \(p\), and
retain its finite
parameter set in \(C\).
Their union \(A'_1\)
is countable. Take a
countable elementary
\(A_1\subseteq B\)
containing \(A_0\cup A'_1\).
Repeat: retain in
\(A'_{m+1}\subseteq C\)
all prior retained
parameters and finite
distinguishing parameters
for each \(a\in A_m\),
and take \(A_{m+1}\)
containing \(A_m\cup
A'_{m+1}\).

Let \(A=\bigcup_mA_m\)
and \(A'=\bigcup_mA'_m\).
The first is a countable
elementary submodel of
\(B\), and \(A'\subseteq
A\cap C\). Restrict
\(p\) to \(A'\) to
obtain \(q\). Every
\(a\in A\) occurs in
some \(A_m\), and its
retained distinguishing
formula prevents it
from realizing \(q\).
The set \(Y\) is
disjoint from \(A'\)
and remains indiscernible
over it, since it was
indiscernible over
the larger \(C\).
This proves Lemma 5.3.
\end{proof}
\end{annexitem}

\begin{annexitem}{annex:p3.morley.omission-transfer}{thm:p03-morley-omission-transfer}{The hard transfer from a countable omission witness}
\begin{proof}
Take \(A,A',Y,q\) from
the preceding lemma.
For the prescribed
\(\kappa>\aleph_0\),
compactness gives a
system containing
\(A'\) and \(\kappa\)
distinct new elements
\(\{y_\beta:\beta<\kappa\}\)
with the same indiscernible
set diagram over \(A'\)
as \(Y\). Every finite
part is realized by
distinct points of
\(Y\) in the original
\(A\). The full diagram
also asserts inequalities
from all elements of
\(A'\), so the new
elements are disjoint
from that subsystem.
Put
\[
 A_\alpha=A'\cup
             \{y_\beta:\beta<\alpha\}
                         \quad(\alpha<\kappa).
\]
This is an actual
continuous chain of
induced parameter
systems. Apply the
continuous-prime-chain
theorem to obtain
\(B_\alpha\) prime
over \(A_\alpha\).

We prove by induction
that every \(B_\alpha\)
omits \(q\). For
finite \(\alpha\),
the finite generator
system embeds over
\(A'\) into \(A\),
using distinct members
of \(Y\). Primality
extends this to an
embedding of \(B_\alpha\)
into \(A\), which
omits \(q\). At a
limit, continuity of
the model chain
preserves omission:
any purported realizing
element already occurs
at an earlier stage.

At an infinite
successor \(\alpha=\beta+1\),
the two generator
sets indexed by
\(\alpha\) and
\(\beta\) have the
same cardinality.
A bijection between
them, fixing \(A'\),
is an isomorphism
\(A_\alpha\to A_\beta\)
because these are
unordered indiscernible
sets over \(A'\).
Primality embeds
\(B_\alpha\) into
\(B_\beta\) extending
that isomorphism.
If \(q\) were realized
in \(B_\alpha\), its
image would realize
it in \(B_\beta\),
contrary to induction.
This successor argument
is exactly where
unordered indiscernibles
and prime-over-the-given-base
are needed.

The union
\(B_\kappa=\bigcup_{\alpha<\kappa}B_\alpha\)
is a model omitting
\(q\). It contains
all \(\kappa\)
distinct generators,
so has size at least
\(\kappa\). In the
prime-chain construction
each sweep adds at
most the size of its
current base plus
\(\aleph_0\); here
successor stages have
size at most
\(\kappa(\alpha)+\aleph_0\),
by induction, and
limits retain this
bound. Hence the
union has size at
most \(\kappa\).
Its omitted type
has the countable
base \(A'\), giving
the required failure
of countable-subsystem
saturation. This is
Theorem 5.4 in full.
\end{proof}
\end{annexitem}

\begin{annexitem}{annex:p3.morley.original-categoricity}{thm:p03-morley-original-categoricity}{The comparison contradiction and final cardinal transfer}
\begin{proof}
Let the categorical
cardinal be
\(\kappa>\aleph_0\).
The thin-model
argument first
makes the theory
totally transcendental.
The realization-chain
construction gives
a model of size
\(\kappa\) saturated
over countable subsystems.
If any uncountable
model were not
fully saturated,
the omission-transfer
theorem would produce
a second model of
size \(\kappa\)
omitting a type
over a countable
subsystem. The two
models cannot be
isomorphic, since
an isomorphism
transports the
parameter subsystem
and all its types.
This contradicts
categoricity. Thus
every uncountable
model is fully
saturated, proving
Theorem 5.5.

For any other
uncountable cardinal
\(\lambda\), models
exist by
Löwenheim--Skolem.
They are all
saturated by the
conclusion just
proved, and any
two of size
\(\lambda\) are
isomorphic by
saturated-model
back-and-forth.
This is Theorem 5.6.
Normalization
transfers both
claims back from
\(\Sigma^*\) to
the original
countable theory.
\end{proof}
\end{annexitem}

\begin{annexitem}{annex:p3.morley.types}{thm:p03-morley-types}{Stone spaces, simultaneous realization, and maps}
\begin{proof}
The quotient formula algebra is Boolean. Its ultrafilters are
exactly maximally consistent one-variable formula sets: the
ultrafilter decision of each Boolean element is the decision of
each formula, and finite meets are nonzero exactly when jointly
consistent. The clopens \(U_\varphi\) give a Hausdorff topology;
compactness is first-order compactness applied to a family of
clopens with the finite-intersection property. There are at most
\(\kappa(A)+\aleph_0\) parameter formulas. Each involves
finitely many named parameters, which proves the finite-base claim.

Naming a new element and asserting its type gives a consistent
extension of \(T(A)\), so compactness realizes each point. To
realize all points, well-order them. At a successor amalgamate a
model realizing the next point with the model already constructed,
over \(A\); at a limit take the elementary union. The resulting
model realizes every point. Löwenheim--Skolem permits the size
bound \(\kappa(A)+\kappa(S(A))+\aleph_0\) when needed.

Parameter substitution is injective on the quotient algebra because
\(f\) is elementary on the parameter systems. Its dual satisfies
\((f^*)^{-1}(U_\varphi)=U_{f\varphi}\), giving continuity.
For surjectivity, amalgamate a realization of any \(p\in S(A)\)
with a model containing \(B\), along \(f\); the realizing element's
type over \(B\) restricts to \(p\). Finally, equality of the
one-types is precisely preservation of every original formula
on the enlarged parameter systems. In the Morleyized language
that is precisely the monomorphism criterion. This proves
Lemma 2.1 and the maps described before Definition 2.2.
\end{proof}
\end{annexitem}

\begin{annexitem}{annex:p3.morley.derivative}{thm:p03-morley-derivative}{Exact images of the global derivatives}
\begin{proof}
The identity map shows that \(\operatorname{Tr}^\beta(A)\)
consists of isolated points of \(S^\beta(A)\); it is thus open
there. Successor derivatives are closed by subtraction of that
open set, and limit derivatives are intersections of closed sets.
Transfinite induction proves closedness and compactness.

We first record a topological fact. If \(h:G\to H\) is a
continuous surjection, \(G\) is compact and \(H\) Hausdorff,
and \(r\) is nonisolated in \(H\), then some point of
\(h^{-1}(r)\) is nonisolated in \(G\). Otherwise this fiber,
being a union of isolated singleton open sets, would be open.
Its complement would be compact, making \(H\setminus\{r\}\)
closed, contrary to the nonisolation of \(r\).

We prove \(f^*(S^\alpha(B))=S^\alpha(A)\) by induction on
\(\alpha\). If \(q\in S^\alpha(B)\) and \(\beta<\alpha\),
then \(q\notin\operatorname{Tr}^\beta(B)\). By the induction
hypothesis \(q\in S^\beta(B)\), so some extension map
\(g:B\to C\) has a nonisolated point \(r\in S^\beta(C)\)
over \(q\). The composite map has that same nonisolated point
over \(f^*(q)\), excluding the latter from
\(\operatorname{Tr}^\beta(A)\). This holds for every
\(\beta<\alpha\), proving inclusion in \(S^\alpha(A)\).

For the reverse inclusion take \(p\in S^\alpha(A)\), and
suppose its fiber in \(S^\alpha(B)\) were empty. The decreasing
compact intersections
\((f^*)^{-1}(p)\cap S^\xi(B)\) start nonempty and become empty.
Their first empty stage cannot be a limit, by compactness.
Thus some \(\beta<\alpha\) is their last nonempty stage, and
\[
 \varnothing\ne(f^*)^{-1}(p)\cap S^\beta(B)
                  \subseteq\operatorname{Tr}^\beta(B).
\]
But \(p\notin\operatorname{Tr}^\beta(A)\); some
\(g:A\to C\) has a nonisolated point \(r\in S^\beta(C)\)
over \(p\). Amalgamate \(B\) and \(C\) over \(A\), obtaining
\(h_1:B\to D\), \(h_2:C\to D\) with \(h_1f=h_2g\).
By the induction hypothesis \(h_2^*\) maps
\(S^\beta(D)\) onto \(S^\beta(C)\). The topological fact
lifts \(r\) to a nonisolated \(s\in S^\beta(D)\).
Then \(h_1^*(s)\) is in the displayed fiber, but cannot be in
\(\operatorname{Tr}^\beta(B)\), since its extension fiber
contains the nonisolated \(s\). This contradiction proves
surjectivity at stage \(\alpha\).

Now use \(\operatorname{Tr}^\alpha=S^\alpha\setminus
S^{\alpha+1}\). Exact images at stages \(\alpha\) and
\(\alpha+1\) show that \(p\) is removed at stage \(\alpha\)
exactly when every point in its stage-\(\alpha\) fiber is
removed there. This is the claimed equivalence.

If \(p\in\operatorname{Tr}^\alpha(A)\), choose a basic clopen
\(U\) with \(U\cap S^\alpha(A)=\{p\}\). It comes from a
finite subsystem \(E\) and a clopen \(V\subseteq S(E)\).
Exact images give \(V\cap S^\alpha(E)=\{i_{EA}^*(p)\}\).
The entire stage-\(\alpha\) fiber of this point in \(S(A)\)
is \(\{p\}\), so the removed-point equivalence puts its
restriction in \(\operatorname{Tr}^\alpha(E)\).
This proves Theorem 2.3 and Corollary 2.4.
\end{proof}
\end{annexitem}

\begin{annexitem}{annex:p3.morley.degree}{thm:p03-morley-degree}{The uniform degree bound and addition}
\begin{proof}
Every stage-\(\alpha\) fiber over a rank-\(\alpha\) point is
compact and consists entirely of isolated points of
\(S^\alpha(B)\); hence it is finite and nonempty. Suppose its
sizes had no common finite bound. For each \(n\) choose an
extension with more than \(n\) such points. Repeated
amalgamation produces an increasing chain of extensions in
which the corresponding fibers have more than \(n\) points.
Exact images ensure that passing to a further extension does
not reduce these cardinalities. In the union the compact fiber
is infinite and therefore contains a nonisolated point, contrary
to the original rank-\(\alpha\) condition. The finite bound
has a largest realized value, \(d(p)\).

For a fixed \(f:A\to B\), there are finitely many
rank-\(\alpha\) points \(q\) over \(p\). Choose an extension
attaining \(d(q)\) for each, and one attaining \(d(p)\).
Finitely many amalgamations put all these extensions in a
common extension \(C\) of \(B\). Each previously attained
fiber size is preserved by exact images, and cannot increase
beyond its maximal degree. Thus the fiber over \(p\) in
\(S^\alpha(C)\) has exactly \(d(p)\) elements, and is the
disjoint union of the fibers over its \(q\)'s, each of size
\(d(q)\). This proves the sum law. All terms are positive,
so a same-rank extension of equal degree leaves no room for
another same-rank point. This is Theorem 2.5, including the
uniformity needed later at chain limits.
\end{proof}
\end{annexitem}

\begin{annexitem}{annex:p3.morley.stabilization}{thm:p03-morley-stabilization}{The two global rank bounds}
\begin{proof}
There are at most \(2^{\aleph_0}\) isomorphism types of finite
parameter systems in the countable relational language, and at
most \(2^{\aleph_0}\) types over each. By finite localization,
every nonempty removal stage occurs among these finite-system
types. Different removal ordinals account for disjoint such
pairs, so there are at most \(2^{\aleph_0}\) nonempty stages.
Moreover an empty global removal stage is permanent: at that
stage every \(S^\alpha(A)\) is unchanged, so the next
definition tests exactly the same spaces and maps. Induction
makes all later removals empty. There cannot therefore be
\((2^{\aleph_0})^+\) consecutive nonempty stages; the first
global stabilization ordinal is below this bound.

If \(S^\alpha(A)=\varnothing\), jointly embed \(A\) and
any \(B\) into \(C\). Exact images first force
\(S^\alpha(C)=\varnothing\), and then
\(S^\alpha(B)=\varnothing\). The least empty stage cannot
be a limit, since a nested family of nonempty closed subsets
of the compact nonempty \(S(A)\) has nonempty intersection.

For the stronger bound in the totally transcendental case,
use the proved countability characterization below. Induct on
\(n\) to show that there are only countably many isomorphism
types of \(n\)-element parameter systems. For \(n=0\)
there is one; each \((n+1)\)-element system is obtained by
adjoining a realization of one of the countably many types
over an \(n\)-element representative. There are thus only
countably many finite-system/type pairs. Every rank localizes
to one of them. Before stabilization every stage is nonempty,
so only countably many ordinals precede \(\alpha_T\).
Consequently \(\alpha_T<\omega_1\). The first argument is
Lemma 2.6; the second is the concluding Theorem 5.7.
\end{proof}
\end{annexitem}

\begin{annexitem}{annex:p3.morley.countability}{thm:p03-morley-countability}{Countable type spaces and the binary residual tree}
\begin{proof}
In the totally transcendental case each point \(p\) has an
ordinal rank \(\alpha\) and a basic clopen \(U(p)\) with
\(U(p)\cap S^\alpha(A)=\{p\}\). These chosen clopens are
distinct. If two points have the same rank this is immediate;
if \(p\) has lower rank than \(q\), then \(q\in S^\alpha(A)\),
so a clopen isolating \(p\) there cannot contain \(q\).
This injects \(S(A)\) into a basis of size at most
\(\kappa(A)+\aleph_0\). Equality modulo \(\aleph_0\)
follows because each named equality \(v_0=a\) gives a
different point. This proves Theorem 2.7.

Conversely suppose the stabilized residuals
\(P(A)=S^{\alpha_T}(A)\) are nonempty. None of their points
is removed. For any residual point there is an extension
fiber containing a nonisolated residual point, in particular
at least two points. A clopen separating two such points
uses finitely many parameters. Exact images thus yield a
finite subsystem whose residual has two disjoint nonempty
clopen pieces.

Construct finite systems \(A_1\subseteq A_2\subseteq\cdots\)
and residual clopens \(U_s\), indexed by binary strings of
length equal to the stage, such that each parent has two
nonempty children in the next extension. Here is why every
parent can be split. Choose a point in that nonempty residual
clopen. Its failure to be removed gives an extension with
two residual points over it; separate them by a clopen and
its complement inside the inverse image of the parent.
At each stage there are only finitely many parents.
Amalgamate their splitting extensions, keeping all old
parameters, and retain the finite parameters of their
separating formulas. Exact images preserve nonemptiness.
For a partition rather than merely two subpieces, assign
the unused part of a parent to one child. This produces
precisely the source's refining binary partitions.

The union \(A\) is countable. For each
\(\eta\in2^\omega\), the inverse images of the clopens
along its branch have the finite-intersection property in
the compact residual \(P(A)\). They have a common point.
Different branches have disjoint intersections, yielding
at least \(2^{\aleph_0}\) points. The space has a countable
basis and hence at most \(2^{\aleph_0}\) points, so equality
holds. This proves Theorem 2.8.

Each parameter set in an original model is an induced
subsystem in the Morleyized language, and translation gives
exactly the same types. If the convention for
\(\omega\)-stability tests only countable models, any
countable parameter set is contained in a countable
elementary submodel; onto restriction transfers the type
bound to that set. The two formulations are thus equivalent.
\end{proof}
\end{annexitem}

\begin{annexitem}{annex:p3.morley.rank-bridge}{thm:p03-morley-rank-bridge}{The splitting rank and the type-space rank agree}
\begin{proof}
For every clopen \(U\subseteq S(A)\), we prove by induction
on \(\alpha\) that its definable set has splitting rank at
least \(\alpha\) exactly when \(U\cap S^\alpha(A)\ne\varnothing\).
The zero case is consistency. At a limit this is compactness
of the nested closed sets \(U\cap S^\beta(A)\).

Suppose \(p\in U\cap S^{\alpha+1}(A)\). Its failure to be
removed at stage \(\alpha\) gives an extension with a
nonisolated \(r\in S^\alpha(B)\) over \(p\). Inside the
inverse image of \(U\), choose infinitely many disjoint
clopens each meeting \(S^\alpha(B)\): recursively separate
a point different from \(r\), retain a clopen neighborhood
of \(r\), and repeat, using nonisolation. By induction the
corresponding definable pieces have rank at least
\(\alpha\), proving the successor inequality.

In the reverse direction, put the parameters of infinitely
many disjoint rank-at-least-\(\alpha\) pieces into one
extension \(B\), using an ambient model or amalgamation.
Their clopens contain distinct points of
\(S^\alpha(B)\). This infinite set has a nonisolated
accumulation point \(r\) in the compact inverse image of
\(U\). Its image is not in \(\operatorname{Tr}^\alpha(A)\),
by the definition testing all extension fibers, and is in
\(S^\alpha(A)\) by exact images. Thus it belongs to
\(U\cap S^{\alpha+1}(A)\).

Under total transcendence, compactness makes the largest
rank on any nonempty clopen exist: if no maximum existed,
the nested intersections up to the supremum would give a
point at that supremum. At the largest rank \(\alpha\),
the clopen intersection with \(S^\alpha(A)\) is compact
and consists of isolated points, hence has finitely many
points \(p_1,\ldots,p_t\). In every extension its
rank-\(\alpha\) points have total multiplicity
\(\sum_jd(p_j)\), by the degree law. Any disjoint
top-rank definable pieces must contain different such
points, bounding their number by this sum. Amalgamate
extensions attaining each \(d(p_j)\); then separate
the resulting finitely many top-rank points by disjoint
clopens. This attains the sum and proves the degree formula.

For finite tuples replace \(S(A)\) throughout by the
space of complete types in a fixed finite tuple of variables.
All Boolean-space, restriction, amalgamation, derivative,
and splitting arguments are unchanged. Countability of
one-types over countable sets implies countability of
\(n\)-types by induction: choose a realization of each
of the countably many possible prefixes; the next
coordinate has only countably many types over that
prefix together with the old countable parameter set.
Two realizations of a prefix can be identified in an
amalgam, so this enumeration covers all continuations.
The residual-tree argument therefore makes every finite
tuple derivative vanish. Conversely ordinal splitting
rank of all unary definable sets precludes a nonempty
residual: a residual point gives a clopen meeting every
derivative and hence a definable set of rank \(\infty\).
The countability characterization then gives
\(\omega\)-stability.

Finally a finite definable set of \(e\) elements has
rank zero: every extension preserves that finite bound,
and its points have uniformly finite isolated fibers.
Conversely a rank-zero clopen has finite top multiplicity
\(e\). If it had more than \(e\) solutions in an
extension, naming them would create more than \(e\)
rank-zero points, a contradiction. Naming all its
solutions also shows its degree equals their number.
\end{proof}
\end{annexitem}

\begin{annexitem}{annex:p3.morley.field-example}{thm:p03-morley-field-example}{The original field illustration and its fixed-characteristic extension}
\begin{proof}
Let \(K=K(A)\). Every algebraic element has an irreducible
minimal polynomial over \(K\). Its distinct roots in any
algebraically closed extension are conjugate over \(K\);
that field isomorphism extends to its algebraic closure
and, after choosing transcendence bases, to an automorphism
of a common algebraically closed extension. Likewise two
elements transcendental over \(K\) have the same type:
the isomorphism \(K(x)\to K(y)\) extends in the same way.
Thus the polynomial equation isolates exactly one type,
and the nonvanishing polynomial conditions specify the
unique remaining type. Rational coefficients in \(K\)
are expressed with finitely many parameters from \(A\)
by clearing their nonzero denominators.

The irreducible polynomial type has finite solution set,
so rank zero. On naming its distinct roots its fiber has
one point per root; no extension can have more roots.
Its degree is precisely the number of distinct roots.
In characteristic zero irreducible polynomials are
separable. In characteristic \(p\), an inseparable
polynomial need not have as many distinct roots as its
polynomial degree, explaining the stated qualification.

There are infinitely many algebraic polynomial types over
every \(K\). They are isolated, and their only possible
accumulation point is the transcendental type. Every
extension has again exactly one nonalgebraic type.
Consequently this type is removed at rank one with degree
one, and the second derivative is empty.

An algebraically closed field is determined up to
isomorphism by its characteristic and transcendence
degree over its prime field. If the dimension is
\(\lambda\), the rational function field on a basis
and its algebraic closure have cardinality
\(\max(\lambda,\aleph_0)\). For uncountable field
cardinality \(\kappa\), this forces dimension
\(\kappa\). For dimension finite or countably infinite,
the field is countably infinite, and the distinct
dimensions give nonisomorphic models. This supplies the
source's Steinitz illustration without confusing
dimension with unary rank.
\end{proof}
\end{annexitem}

\begin{annexitem}{annex:p3.morley.tuple-code-example}{thm:p03-morley-tuple-code-example}{Coordinate-counting ranks in Example II}
\begin{proof}
Every bijection of the infinite \(R_0\)-part extends
uniquely to the codes, and every isomorphism restricts
to such a bijection. Its finite Cartesian power has the
same infinite cardinality, proving categoricity in
every infinite size. The elementary theory is complete
because finite tuples are described by finite equality,
coordinate, and code patterns; a new unused base point
always remains available.

Taking coordinates and forming codes determines unique
elements from the named parameters, so every element
of \(\overline B\) is individually definable over
\(B\) and has rank zero. Any permutation of the unused
base points extends to an automorphism fixing
\(\overline B\). Thus a fresh base point has one type,
and a code's type is determined by its slots occupied
by points of \(\overline B\), their values, and the
distinctness of the remaining coordinates.

Induct on the number \(k\) of fresh coordinates.
The zero case is algebraic. For \(k\geq1\), naming
one of infinitely many possible new values in a fixed
fresh slot gives infinitely many disjoint extension
pieces with \(k-1\) fresh coordinates. The splitting
criterion gives rank at least \(k\). In every
extension the type either acquires some named
coordinates, then has rank less than \(k\) by
induction, or retains all \(k\) coordinates fresh,
then has the single generic type of that slot pattern.
This latter point is isolated among the stage-\(k\)
points in the clopen specifying the old slot pattern;
therefore the original point is removed at stage
\(k\), giving exact rank \(k\). The same argument
for a fresh base point gives exact rank one.
The maximum rank is \(n\), so the first empty stage
is \(n+1\). The assumption \(n\geq1\) is material:
with zero-coordinate codes the infinite base still
has rank one, not rank zero.
\end{proof}
\end{annexitem}

\begin{annexitem}{annex:p3.morley.unary-example}{thm:p03-morley-unary-example}{Example III and ordinal addition}
\begin{proof}
Finite formulas in this unary language reduce to Boolean
combinations of predicate cells and equalities. Quantifiers
are eliminated by noting that every allowed cell has
infinitely many elements, so excluding finitely many
named elements never exhausts one. This proves completeness
and the description of one-types. Named elements give
isolated algebraic points. Each other point is determined
by a predicate sequence in \(Y\) and inequalities from
all named elements. A sequence outside \(Y\) has a
finite forbidden cylinder, so is inconsistent; each
sequence in \(Y\) is consistent by compactness.

No fresh-element point is algebraic. It has arbitrarily
many distinct realizations in an extension; naming
infinitely many produces a nonisolated fresh point in
the rank-zero fiber. Conversely every finite definable
set consists of named elements. Thus the first
derivative is exactly the fresh-point space. Its
predicate-cylinder topology is \(Y\)'s topology.
Restriction between fresh-point spaces forgets no
predicate information and is a homeomorphism. The
later global derivative therefore removes exactly
the ordinarily isolated points of the corresponding
derivative of \(Y\), proving
\(S^{1+\beta}(A)\cong Y^{(\beta)}\).

For completeness, the familiar countability assertion
for closed subsets of Cantor space follows from its
countable basis. Each removed point at any stage has
a basic clopen isolating it at that stage. A basic
clopen cannot be chosen for different removed points,
by the same rank comparison as in the countability
proof. Hence only countably many points are removed.
A nonempty perfect residual contains a binary tree
of separated nested clopens, and compactness gives
a different point for each infinite branch. It is
therefore uncountable. The residual is empty exactly
when \(Y\) is countable. Stabilization of the two
derivatives now gives \(\alpha_T=1+\alpha_Y\).
\end{proof}
\end{annexitem}
